% ===========================================================================
%  Raport — Faza 2: Cuantizare INT8 Weight-Only (per-tensor)
%  Evaluare ViT-Tiny pretrained pe ImageNet-1k validation
%  Tudose Alexandru · Lucrare de Licență 2026
% ===========================================================================
\documentclass[12pt,a4paper]{article}

% --- Encoding & fonts (XeTeX / tectonic) ---
\usepackage{fontspec}
\defaultfontfeatures{Ligatures=TeX}

% --- Language ---
\usepackage{polyglossia}
\setmainlanguage{romanian}

% --- Page geometry ---
\usepackage[
  top=2.5cm, bottom=2.5cm,
  left=2.5cm, right=2.5cm
]{geometry}

% --- Math ---
\usepackage{amsmath, amssymb}

% --- Graphics ---
\usepackage{graphicx}
\usepackage{float}
\usepackage{caption}
\usepackage{subcaption}
\captionsetup{font=small, labelfont=bf}

% --- Tables ---
\usepackage{booktabs}
\usepackage{array}
\usepackage{multirow}

% --- Colors ---
\usepackage{xcolor}
\definecolor{fp32color}{HTML}{0072B2}
\definecolor{fp16color}{HTML}{009E73}
\definecolor{int8color}{HTML}{E69F00}
\definecolor{lightgray}{HTML}{F4F4F4}

% --- Hyperlinks ---
\usepackage[
  colorlinks=true,
  linkcolor=fp32color,
  citecolor=fp32color,
  urlcolor=fp32color
]{hyperref}

% --- Misc ---
\usepackage{microtype}
\usepackage{parskip}
\setlength{\parindent}{0pt}
\setlength{\parskip}{6pt}
\usepackage{enumitem}

% --- Header / Footer ---
\usepackage{fancyhdr}
\pagestyle{fancy}
\fancyhf{}
\lhead{\small Tudose Alexandru --- Faza 2: INT8 Weight-Only Quantization}
\rhead{\small 2026}
\cfoot{\thepage}
\renewcommand{\headrulewidth}{0.4pt}

% ===========================================================================
\begin{document}
% ===========================================================================

% --- Title page ---
\begin{titlepage}
  \centering
  {\large Universitatea Transilvania din Brașov\\}
  {\small Facultatea de Matematică și Informatică\\}
  {\small Specializarea: Informatică Aplicată\\[2cm]}

  \rule{\linewidth}{0.4pt}\\[0.5cm]
  {\LARGE\bfseries Faza 2: Cuantizare INT8 Weight-Only\\[0.3cm]}
  {\Large\bfseries Per-Tensor cu Scalare Liniară\\[0.2cm]}
  {\Large\bfseries ViT-Tiny pe ImageNet-1k validation\\[0.2cm]}
  {\large Post-Training Quantization --- FP32 vs.\ INT8}\\[0.3cm]
  \rule{\linewidth}{0.4pt}\\[2cm]

  \begin{minipage}{0.45\textwidth}
    \textbf{Student:}\\
    Tudose Alexandru\\
    Grupa 10LF333\\
    An III, Informatică Aplicată
  \end{minipage}
  \hfill
  \begin{minipage}{0.45\textwidth}
    \raggedleft
    \textbf{Coordonator:}\\
    Conf.\ dr.\ ing.\ Honorius Gâlmeanu\\
    Informatică Aplicată
  \end{minipage}\\[3cm]

\end{titlepage}

% --- TOC ---
\tableofcontents
\newpage

% ===========================================================================
\section{Introducere}
% ===========================================================================

Acest raport documentează \textbf{Faza~2} a lucrării de licență, în care
evaluăm cuantizarea \textbf{INT8 weight-only} cu scalare liniară per-tensor
pe ViT-Tiny pretrained. Spre deosebire de Faza~1 (FP16), care reduce
precizia la 16 biți, INT8 comprimă ponderile la doar \textbf{8 biți},
oferind o reducere teoretică de $4\times$ a memoriei.

\textbf{Obiective:}
\begin{enumerate}
  \item Implementarea cuantizării INT8 cu scalare liniară per-tensor
    pentru straturile \texttt{nn.Linear} ale ViT-Tiny.
  \item Evaluarea impactului asupra acurateții, latenței și memoriei
    față de baseline-ul FP32 stabilit în Faza~1.
  \item Analiza erorii de cuantizare per strat (MSE, MAE, scale factor)
    pentru identificarea straturilor problematice.
  \item Comparația cross-phase FP32 / FP16 / INT8.
\end{enumerate}

\textbf{Context:} Cuantizarea weight-only INT8 stochează ponderile ca
\texttt{int8} (1 byte/param) dar efectuează inferența în \texttt{float32}
prin dequantizare la fiecare forward pass. Aceasta reduce memoria, dar
\emph{nu} accelerează inferența --- overhead-ul de dequantizare poate chiar
încetini modelul.

% ===========================================================================
\section{Metodologie}
% ===========================================================================

\subsection{Schema de cuantizare}

Cuantizarea utilizează scalare liniară per-tensor, cu formulele:

\begin{equation}
  \text{scale} = \frac{\max(|W|)}{127}
  \label{eq:scale}
\end{equation}

\begin{equation}
  q = \text{round}\!\left(\frac{W}{\text{scale}}\right)
    .\text{clamp}(-128, 127).\text{to}(\texttt{int8})
  \label{eq:quantize}
\end{equation}

\begin{equation}
  \hat{W} = q.\text{float()} \times \text{scale}
  \label{eq:dequantize}
\end{equation}

Ecuația~\eqref{eq:dequantize} este aplicată la fiecare forward pass:
ponderile INT8 sunt dequantizate la \texttt{float32} înainte de multiplicarea
matriceală. Aceasta garantează o mapare bijectivă la 256 valori discrete,
spre deosebire de formatul \texttt{float8\_e4m3fn} testat în Faza~0.

\subsection{Cuantizare selectivă}

Nu toate straturile sunt cuantizate. Straturile sensibile la erori numerice
sunt excluse:

\begin{table}[H]
  \centering
  \caption{Straturi cuantizate vs.\ excluse}
  \begin{tabular}{lll}
    \toprule
    \textbf{Categorie} & \textbf{Straturi} & \textbf{Motiv} \\
    \midrule
    \textcolor{int8color}{\textbf{Cuantizate}} &
      \texttt{attn.qkv}, \texttt{attn.proj} &
      Proiecții liniare --- robuste la cuantizare \\
    & \texttt{mlp.fc1}, \texttt{mlp.fc2} &
      MLP layers --- cea mai mare parte din parametri \\
    \midrule
    \textcolor{fp32color}{\textbf{Excluse (FP32)}} &
      \texttt{norm1}, \texttt{norm2}, \texttt{norm} &
      LayerNorm --- sensibil la erori de scalare \\
    & \texttt{cls\_token}, \texttt{pos\_embed} &
      Parametri de embedding --- critici pentru clasificare \\
    & \texttt{head} &
      Classification head --- stratul final de decizie \\
    \bottomrule
  \end{tabular}
  \label{tab:skip_patterns}
\end{table}

\textbf{Rezultat:} 48 din 49 straturi \texttt{nn.Linear} sunt cuantizate
(5\,308\,416 parametri). Doar \texttt{head} (192\,000 parametri) rămâne
în FP32.

\subsection{Configurația experimentală}

\begin{table}[H]
  \centering
  \caption{Configurația experimentală}
  \begin{tabular}{ll}
    \toprule
    \textbf{Parametru} & \textbf{Valoare} \\
    \midrule
    Model             & \texttt{vit\_tiny\_patch16\_224} (timm) \\
    Parametri         & 5.7 milioane (din care 5.3M cuantizate) \\
    Pretrained        & ImageNet-1k (\texttt{pretrained=True}) \\
    Dataset evaluare  & ImageNet-1k validation (50\,000 imagini, 1\,000 clase) \\
    Batch size        & 64 \\
    Device            & Apple M4, backend MPS \\
    PyTorch           & 2.9.1 \\
    Metoda cuantizare & Per-tensor INT8, scalare liniară \\
    Granularitate     & Un factor de scalare per matrice de ponderi \\
    \bottomrule
  \end{tabular}
  \label{tab:config_int8}
\end{table}

\subsection{Metrici de evaluare}

Pe lângă metricile standard (accuracy, loss, latency, memory), Faza~2
introduce metrici specifice cuantizării:

\begin{itemize}
  \item \textbf{MSE} (Mean Squared Error): $\frac{1}{n}\sum(W - \hat{W})^2$
    --- eroarea medie de reconstrucție per strat.
  \item \textbf{MAE} (Mean Absolute Error): $\frac{1}{n}\sum|W - \hat{W}|$
    --- eroarea absolută medie.
  \item \textbf{Scale factor}: Valoarea $s = \max(|W|)/127$ per strat ---
    reflectă dinamicul ponderilor.
\end{itemize}

% ===========================================================================
\section{Rezultate}
% ===========================================================================

\subsection{Comparație globală FP32 vs.\ INT8}

\begin{table}[H]
  \centering
  \caption{Rezultate FP32 vs.\ INT8 pe ImageNet-1k validation}
  \begin{tabular}{lrrrr}
    \toprule
    \textbf{Format} & \textbf{Accuracy (\%)} & \textbf{Loss} &
    \textbf{Latency (ms)} & \textbf{Memory (MB)} \\
    \midrule
    FP32 (baseline) & 75.456 & 0.9420 & 200.99 & 21.81 \\
    INT8 (per-tensor) & 75.134 & 0.9563 & 217.09 & 6.62 \\
    \midrule
    \textbf{Diferență} & $\mathbf{-0.322}$~pp & $+0.014$ &
    $\mathbf{0.93\times}$ (mai lent) & $\mathbf{3.29\times}$ reducere \\
    \bottomrule
  \end{tabular}
  \label{tab:int8_results}
\end{table}

\begin{figure}[H]
  \centering
  \includegraphics[width=\textwidth]{../results/INT8ImageNet/plots/00_summary.png}
  \caption{Overview Faza~2: acuratețe, memorie și latență --- FP32 vs.\ INT8
    per-tensor. Degradarea de acuratețe este mică ($-0.322$~pp), reducerea de
    memorie este semnificativă ($3.29\times$), dar nu există speedup
    (overhead de dequantizare).}
  \label{fig:int8_summary}
\end{figure}

\subsection{Acuratețe și loss}

Degradarea de \textbf{$-0.322$~pp} (de la 75.456\% la 75.134\%) este
acceptabilă pentru o reducere de memorie de $3.29\times$. Aceasta
reprezintă 161 imagini din 50\,000 clasificate diferit.

Loss-ul crește de la 0.9420 la 0.9563 ($+1.5\%$), indicând o ușoară
``estompare'' a distribuției softmax cauzată de eroarea de cuantizare.
Totuși, această creștere nu se traduce într-o degradare practică
semnificativă a acurateții.

\subsection{Memorie}

Reducerea de $\mathbf{3.29\times}$ (de la 21.81~MB la 6.62~MB) este
explicată astfel:

\begin{itemize}
  \item 48 straturi cuantizate: $5\,308\,416 \times 1$ byte $= 5.06$~MB
    (față de $5\,308\,416 \times 4$ bytes $= 20.25$~MB în FP32)
  \item Straturi excluse (head, bias-uri): rămân în FP32
  \item Buffer-uri suplimentare: scale factors (48 valori float32) ---
    overhead neglijabil
  \item Reducerea nu este exact $4\times$ deoarece nu toate componentele
    sunt cuantizate
\end{itemize}

\subsection{Latență}

\textbf{INT8 este mai lent} decât FP32 cu un factor de $0.93\times$
(217.09~ms vs.\ 200.99~ms). Aceasta este o consecință directă a
cuantizării weight-only:

\begin{enumerate}
  \item La fiecare forward pass, ponderile \texttt{int8} sunt dequantizate
    la \texttt{float32} ($\hat{W} = q \cdot s$).
  \item Multiplicarea matriceală se execută în \texttt{float32},
    ca la modelul original.
  \item Overhead-ul de dequantizare ($\approx 16$~ms/batch) anulează orice
    beneficiu de bandwidth.
\end{enumerate}

\textbf{Concluzie:} Pentru a obține speedup real cu INT8, ar fi necesară
inferență nativă în aritmetică pe numere întregi (e.g., NVIDIA INT8 Tensor
Cores sau \texttt{torch.int8} GEMM kernels), care nu sunt disponibile pe
Apple M4 MPS.

\subsection{Eroare de cuantizare per strat}

\begin{table}[H]
  \centering
  \caption{Statistici sumare ale erorii de cuantizare}
  \begin{tabular}{lr}
    \toprule
    \textbf{Metrică} & \textbf{Valoare} \\
    \midrule
    Straturi cuantizate         & 48 \\
    Parametri cuantizați        & 5\,308\,416 \\
    MSE mediu per strat         & $2.88 \times 10^{-6}$ \\
    Worst layer (MSE)           & \texttt{blocks.7.mlp.fc2} ($3.54 \times 10^{-5}$) \\
    Best layer (MSE)            & \texttt{blocks.8.attn.proj} ($2.69 \times 10^{-7}$) \\
    Raport worst/best           & $131.6\times$ \\
    \bottomrule
  \end{tabular}
  \label{tab:quant_stats}
\end{table}

\begin{figure}[H]
  \centering
  \includegraphics[width=\textwidth]{../results/INT8ImageNet/plots/01_per_layer_mse.png}
  \caption{MSE per strat cuantizat. Straturile \texttt{blocks.7.mlp.fc1} și
    \texttt{blocks.7.mlp.fc2} sunt outlieri evidenți (MSE $\approx 12\times$
    mai mare decât media), marcate cu portocaliu (>P90). Aceste straturi au
    scale factors de $\approx 0.020$, indicând ponderi cu valori absolute
    mult mai mari decât restul modelului.}
  \label{fig:per_layer_mse}
\end{figure}

\subsection{Distribuția scale factors}

\begin{figure}[H]
  \centering
  \includegraphics[width=\textwidth]{../results/INT8ImageNet/plots/02_per_layer_scale.png}
  \caption{Distribuția scale factors per strat. (Stânga) Valori per strat
    în ordinea modelului --- se observă spike-urile la block~7. (Dreapta)
    Boxplot per tip de strat: \texttt{attn.qkv} și \texttt{mlp.fc1/fc2} au
    variabilitate mai mare decât \texttt{attn.proj}.}
  \label{fig:scale_distribution}
\end{figure}

Outlierii din block~7 au scale factors de $\approx 0.020$, comparativ cu
media de $\approx 0.004$. Aceasta înseamnă că ponderile acestor straturi
au valori absolute de $5\times$ mai mari decât media, ceea ce forțează
o ``grilă'' de cuantizare mult mai grosolană (fiecare nivel INT8
acoperă un interval mai larg).

\textbf{Implicație:} Aceste straturi sunt candidați ideali pentru
cuantizare \emph{per-channel} (un scale factor per rând), care va fi
evaluată în Faza~3.

\subsection{Comparație cross-phase: FP32 / FP16 / INT8}

\begin{figure}[H]
  \centering
  \includegraphics[width=\textwidth]{../results/INT8ImageNet/plots/03_cross_phase_comparison.png}
  \caption{Comparație între cele trei formate evaluate până acum. FP16
    oferă cel mai bun raport acuratețe/eficiență (zero degradare, $2\times$
    memorie, $1.22\times$ speedup). INT8 reduce memoria mai agresiv
    ($3.29\times$), dar cu o ușoară degradare și fără speedup.}
  \label{fig:cross_phase}
\end{figure}

\begin{table}[H]
  \centering
  \caption{Comparație cumulativă Faza~1 + Faza~2}
  \begin{tabular}{lrrrr}
    \toprule
    \textbf{Format} & \textbf{Accuracy} & \textbf{Degradare} &
    \textbf{Memorie} & \textbf{Speedup} \\
    \midrule
    FP32            & 75.456\% & ---           & 21.81 MB & $1.00\times$ \\
    FP16 (Faza 1)   & 75.452\% & $-0.004$~pp  & 10.91 MB ($2.00\times$) & $1.22\times$ \\
    INT8 (Faza 2)   & 75.134\% & $-0.322$~pp  & \phantom{0}6.62 MB ($3.29\times$) & $0.93\times$ \\
    \bottomrule
  \end{tabular}
  \label{tab:cross_phase}
\end{table}

\subsection{Tradeoff acuratețe--eficiență}

\begin{figure}[H]
  \centering
  \includegraphics[width=\textwidth]{../results/INT8ImageNet/plots/04_tradeoff.png}
  \caption{Scatter plot acuratețe vs.\ memorie (stânga) și acuratețe vs.\
    latență (dreapta). FP16 domină pe ambele axe. INT8 oferă o reducere de
    memorie suplimentară față de FP16, dar cu un cost în latență.}
  \label{fig:tradeoff}
\end{figure}

% ===========================================================================
\section{Analiza Outlierilor --- Block 7}
% ===========================================================================

Straturile \texttt{blocks.7.mlp.fc1} și \texttt{blocks.7.mlp.fc2}
prezintă un comportament atipic:

\begin{table}[H]
  \centering
  \caption{Outlieri vs.\ media modelului}
  \begin{tabular}{lrrrr}
    \toprule
    \textbf{Strat} & \textbf{MSE} & \textbf{MAE} &
    \textbf{Scale} & \textbf{Raport MSE/medie} \\
    \midrule
    Media model & $2.88 \times 10^{-6}$ & --- & 0.004 & $1.0\times$ \\
    \texttt{b7.mlp.fc1} & $3.32 \times 10^{-5}$ & 0.0050 & 0.020 & $11.5\times$ \\
    \texttt{b7.mlp.fc2} & $3.54 \times 10^{-5}$ & 0.0052 & 0.021 & $12.3\times$ \\
    \bottomrule
  \end{tabular}
  \label{tab:block7_outliers}
\end{table}

\textbf{Cauza:} Aceste straturi conțin ponderi cu valori absolute
semnificativ mai mari ($\max|W| \approx 2.5$, comparativ cu $\approx 0.5$
pentru straturile tipice). Cuantizarea per-tensor alocă un singur scale
factor pentru întreaga matrice, ceea ce înseamnă că valorile mici sunt
cuantizate cu o rezoluție foarte scăzută.

\textbf{Soluție:} Cuantizarea \emph{per-channel} (un scale factor per rând
al matricii) va fi evaluată în Faza~3. Aceasta permite fiecărui canal de
ieșire să aibă propriul factor de scalare, adaptându-se la dinamicul local
al ponderilor.

% ===========================================================================
\section{Discuție}
% ===========================================================================

\subsection{Validarea cuantizării selective}

Strategia de a exclude LayerNorm, cls\_token, pos\_embed și head
s-a dovedit corectă: degradarea de doar $-0.322$~pp pentru 48 straturi
cuantizate confirmă că straturile liniare din Attention și MLP sunt
robuste la cuantizare INT8.

\subsection{Weight-only vs.\ activation quantization}

Implementarea noastră cuantizează doar \emph{ponderile}. Activările
rămân în \texttt{float32} pe tot parcursul inferenței. Aceasta:
\begin{itemize}
  \item \textbf{Avantaj:} Simplitate --- nu necesită calibrare pe date
    de referință, nu depinde de distribuția intrărilor.
  \item \textbf{Dezavantaj:} Nu oferă speedup pe niciun hardware,
    deoarece GEMM-urile rămân în float32.
\end{itemize}

Pentru speedup real ar fi necesară cuantizarea \emph{completă} (ponderi
+ activări), cu suport hardware nativ (e.g., NVIDIA \texttt{int8} GEMM
via TensorRT sau ONNX Runtime).

\subsection{Comparație cu Faza~0}

\begin{table}[H]
  \centering
  \caption{Progresul cuantizării: Faza~0 vs.\ Faza~2}
  \begin{tabular}{lll}
    \toprule
    \textbf{Aspect} & \textbf{Faza 0 (CIFAR-10)} & \textbf{Faza 2 (ImageNet-1k)} \\
    \midrule
    Format              & FP8 E4M3FN             & INT8 scalare liniară \\
    Dataset             & 10\,000 imagini, 10 cls & 50\,000 imagini, 1\,000 cls \\
    Cuantizare          & Toate straturile        & Selectivă (skip norm/head) \\
    Mapare              & Non-bijectivă           & Bijectivă (256 valori) \\
    Degradare           & $<0.02$~pp              & $-0.322$~pp \\
    Analiză per-strat   & Nu                      & Da (MSE, MAE, scale) \\
    \bottomrule
  \end{tabular}
  \label{tab:faza0_vs_faza2}
\end{table}

Degradarea mai mare pe ImageNet ($-0.322$~pp vs.\ $<0.02$~pp pe CIFAR-10)
confirmă ipoteza din Faza~0: \textbf{dataset-urile complexe sunt mai
sensibile la erori de cuantizare} decât dataset-urile simple.

\subsection{Limitări}

\begin{enumerate}
  \item \textbf{Granularitate per-tensor:} Un singur scale factor per
    matrice nu se adaptează la distribuții non-uniforme ale ponderilor.
    Faza~3 va evalua alternativa per-channel.
  \item \textbf{Fără speedup:} Dequantizarea la \texttt{float32} la
    fiecare forward pass introduce overhead. Aceasta este o limitare
    fundamentală a cuantizării weight-only pe hardware fără suport INT8
    nativ.
  \item \textbf{Hardware-specific:} Rezultatele de latență sunt specifice
    Apple M4 MPS. Pe GPU-uri NVIDIA cu INT8 Tensor Cores, profilul de
    latență ar fi diferit.
  \item \textbf{Un singur model:} Doar ViT-Tiny a fost evaluat. Modele
    mai mari pot prezenta pattern-uri diferite de distribuție a ponderilor.
\end{enumerate}

% ===========================================================================
\section{Concluzii}
% ===========================================================================

\begin{enumerate}
  \item \textbf{Degradare acceptabilă:} INT8 per-tensor introduce o
    pierdere de $\mathbf{-0.322}$~pp (161 imagini din 50\,000),
    acceptabilă pentru o reducere semnificativă de memorie.

  \item \textbf{Reducere de memorie $\mathbf{3.29\times}$:} De la 21.81~MB
    la 6.62~MB. Nu exact $4\times$ teoretic, deoarece bias-urile și head-ul
    rămân în FP32.

  \item \textbf{Fără speedup pe MPS:} INT8 weight-only este chiar
    $\mathbf{7\%}$ mai lent decât FP32, din cauza overhead-ului de
    dequantizare la fiecare forward pass.

  \item \textbf{Outlieri identificați:} Straturile \texttt{blocks.7.mlp.*}
    au MSE de $\approx 12\times$ media, cu scale factors de $5\times$
    mai mari. Acestea sunt candidații principali pentru cuantizare per-channel.

  \item \textbf{Cuantizare selectivă validată:} Excluderea LayerNorm,
    cls\_token, pos\_embed și head menține acuratețea la un nivel acceptabil,
    confirmând robustețea straturilor liniare din Attention și MLP.

  \item \textbf{Direcții Faza~3:} Analiza de sensibilitate per strat
    (câte un strat cuantizat la un moment dat), comparația per-tensor vs.\
    per-channel, și profilarea detaliată a memoriei și latenței.
\end{enumerate}

% ===========================================================================
\begin{thebibliography}{9}
% ===========================================================================

\bibitem{dosovitskiy2021vit}
A. Dosovitskiy, L. Beyer, A. Kolesnikov et al.,
``An Image is Worth 16x16 Words: Transformers for Image Recognition at Scale,''
\emph{ICLR}, 2021. \href{https://arxiv.org/abs/2010.11929}{arXiv:2010.11929}.

\bibitem{jacob2018quantization}
B. Jacob, S. Kligys, B. Chen et al.,
``Quantization and Training of Neural Networks for Efficient
Integer-Arithmetic-Only Inference,''
\emph{CVPR}, 2018. \href{https://arxiv.org/abs/1712.05877}{arXiv:1712.05877}.

\bibitem{nagel2021whitepaper}
M. Nagel, M. Fournarakis, R. A. Amjad, et al.,
``A White Paper on Neural Network Quantization,''
2021. \href{https://arxiv.org/abs/2106.08295}{arXiv:2106.08295}.

\bibitem{dettmers2022llmint8}
T. Dettmers, M. Lewis, Y. Belkada, L. Zettlemoyer,
``LLM.int8(): 8-bit Matrix Multiplication for Transformers at Scale,''
\emph{NeurIPS}, 2022. \href{https://arxiv.org/abs/2208.07339}{arXiv:2208.07339}.

\bibitem{timm}
R. Wightman, ``PyTorch Image Models (timm),'' 2019.
\url{https://github.com/rwightman/pytorch-image-models}.

\bibitem{imagenet}
O. Russakovsky, J. Deng, H. Su et al.,
``ImageNet Large Scale Visual Recognition Challenge,''
\emph{IJCV}, vol.~115, no.~3, pp.~211--252, 2015.
\href{https://arxiv.org/abs/1409.0575}{arXiv:1409.0575}.

\end{thebibliography}

\end{document}
